\documentclass[../../main.tex]{subfiles}
\begin{document}

\begin{center}
  {\fontsize{48pt}{58pt}\selectfont \textbf{2012}\normalsize\textbf{年度}}
\end{center}

\begin{table}[htb]
  \centering
  \caption{出題テーマと難易度}
  \begin{tabular}{|c|p{9cm}|c|}
    \hline
    \textbf{問題番号} & \textbf{テーマ} & \textbf{難易度} \\
    \hline
    第1問 & 正四面体のベクトルの内積，さいころの確率 & \\
    \hline
    第2問 & 等比数列の桁数，平方根の整数部分と約数 & \\
    \hline
    第3問 & 三次関数と直線の交点，囲まれる面積の最小値 & \\
    \hline
    第4問 & 漸化式で定めた数列と和の極限 & \\
    \hline
    第5問 & 行列による一次変換と距離の比 & \\
    \hline
    第6問 & 四面体の一部分の体積 & \\
    \hline
  \end{tabular}
\end{table}

\begin{table}[htb]
  \centering
  \caption{解答}
  \renewcommand{\arraystretch}{1.4}
  \begin{tabular}{|c|c|p{9cm}|}
    \hline
    \textbf{問題番号} & \textbf{小問} & \textbf{答え} \\
    \hline
    第1問 & (1) & $\overrightarrow{DE}\cdot\overrightarrow{AC}=\dfrac{1}{2}$ \\
    \cline{2-3}
    & (2) & 確率 $\dfrac{1}{3}$ \\
    \hline
    第2問 & (1) & $48$桁 \\
    \cline{2-3}
    & (2) & $298$個 \\
    \hline
    第3問 & (1) & $a\ge-\dfrac{1}{4}$ \\
    \cline{2-3}
    & (2) & $a=38-27\sqrt{2}$ \\
    \hline
    第4問 & (1) & $a_2=\dfrac{1}{(n+1)(n+2)}$，
    $a_3=\dfrac{1}{(n+2)(n+3)}$ \\
    \cline{2-3}
    & (2) & $a_k=\dfrac{1}{(n+k-1)(n+k)}$ \\
    \cline{2-3}
    & (3) & $\displaystyle\lim_{n\to\infty}b_n=\log 2$（証明問題） \\
    \hline
    第5問 & (1)，(2) & \\
    \hline
    第6問 & -- & $V=4\sqrt{3}-2\pi$ \\
    \hline
  \end{tabular}
\end{table}

\paragraph{試験形態}
6問．試験時間は要確認．

\end{document}
